\documentclass[border=12pt]{standalone}\usepackage[siunitx, european]{circuitikz}\usepackage{amsmath}
\begin{document}

% 四个BJT: 仅标注B、C、E对地电位
\begin{circuitikz}[scale=1.0,transform shape]
  \draw (0,0) node[npn](Q){};
  \draw (Q.B) -- ++(-1,0) node[left] {5V};
  \draw (Q.C) -- ++(0.5,0) node[right] {4.4V};
  \draw (Q.E) -- ++(0.5,0) node[right] {4.3V};
  \draw (0.8,-1.2) node[font=\small] {(a)};
\end{circuitikz}
\begin{circuitikz}[scale=1.0,transform shape]
  \draw (0,0) node[pnp](Q){};
  \draw (Q.B) -- ++(-1,0) node[left] {-2V};
  \draw (Q.C) -- ++(0.5,0) node[right] {-7V};
  \draw (Q.E) -- ++(0.5,0) node[right] {-1.6V};
  \draw (0.8,-1.2) node[font=\small] {(b)};
\end{circuitikz}
\begin{circuitikz}[scale=1.0,transform shape]
  \draw (0,0) node[npn](Q){};
  \draw (Q.B) -- ++(-1,0) node[left] {1.3V};
  \draw (Q.C) -- ++(0.5,0) node[right] {4V};
  \draw (Q.E) -- ++(0.5,0) node[right] {2V};
  \draw (0.8,-1.2) node[font=\small] {(c)};
\end{circuitikz}
\begin{circuitikz}[scale=1.0,transform shape]
  \draw (0,0) node[npn](Q){};
  \draw (Q.B) -- ++(-1,0) node[left] {3V};
  \draw (Q.C) -- ++(0.5,0) node[right] {8.3V};
  \draw (Q.E) -- ++(0.5,0) node[right] {2.3V};
  \draw (0.8,-1.2) node[font=\small] {(d)};
\end{circuitikz}

\end{document}